% Chapter of the map: Physics: the Ising model.
% Conventions and the template of an entry: blueprint/README.md. Check with
% blueprint/tools/check_names.sh. Labels of this file begin with kind:physics (def:physics-..., thm:physics-...).
\chapter{Physics: the Ising model}
\label{chap:physics}

This chapter collects the statements of the map that read a computational problem as a physical
system. They are of two kinds. Some are about the Ising model without external field: its
ground-state decision problem is $\cNP$-complete, the energy of a configuration is the weight of a
cut up to an affine change, and a structure-preserving bijection carries ground states,
stabilizers and orbit counts from one system to another. The others are about finite systems:
the first and second moments and the annealed bound for random $3$-SAT, the sum over programs of
a finite model as a Gibbs mass, and the recurrence of an injective map on a finite set. Each
section says in which model its statements live; none of them is about a physical system in the
continuum.

\section{The Ising ground-state problem}
\label{sec:physics-ising-problem}

An instance is a list of couplings between sites and a threshold; the question is whether some
configuration of spins has energy at most the threshold. The statements of this section are about
Turing-machine complexity classes, through the definitions of the classes in Part~I.

\begin{definition}[The problem $\ISING$: ground-state energy at most a threshold]
  \label{def:physics-ising-problem}
  \lean{Sagitta.Map.Informacion.Hamiltoniano.Acople, Sagitta.Map.Informacion.Hamiltoniano.IsingInstancia, Sagitta.Map.Informacion.Hamiltoniano.espin, Sagitta.Map.Informacion.Hamiltoniano.energia, Sagitta.Map.Informacion.Hamiltoniano.ISING, Sagitta.Map.Informacion.Hamiltoniano.encIsing}
  \leanok
  A \emph{coupling} is a triple $(u, v, J)$ with $u, v \in \mathbb{N}$ (sites) and
  $J \in \mathbb{Z}$. An instance is a list $E$ of couplings together with a threshold
  $t \in \mathbb{Z}$. A configuration is a map from sites to $\{\mathrm{true}, \mathrm{false}\}$; the
  spin of a site is $+1$ on $\mathrm{true}$ and $-1$ on $\mathrm{false}$. The energy of a
  configuration $\sigma$ is
  \[
    H_E(\sigma) = -\sum_{(u,v,J) \in E} J\, \sigma_u \sigma_v ,
  \]
  with no external field. $\ISING$ is the set of instances for which some finite set $S$ of sites,
  with $\sigma_v = +1$ exactly when $v \in S$, has $H_E(\sigma) \le t$. The encoding of an instance
  writes the list of couplings, each as $(u, v, J)$ in binary (the integer with sign and magnitude),
  followed by $t$ in binary.
  \own{}
  \leanfile{Sagitta/Map/Informacion/Hamiltoniano.lean}
\end{definition}

\begin{theorem}[$\ISING$ is $\cNP$-complete, for the encoding of lists of couplings]
  \label{thm:physics-ising-npcomplete}
  \lean{Sagitta.Map.Informacion.Hamiltoniano.ISING_npComplete}
  \leanok
  \uses{def:physics-ising-problem, thm:maxcut, def:NP, def:reduction}
  The problem $\ISING$ of Definition~\ref{def:physics-ising-problem} is $\cNP$-complete for the
  encoding of that definition (in the sense of the Turing-machine definitions of
  Part~I).
  \own{}
  \leanfile{Sagitta/Map/Informacion/Hamiltoniano.lean}
\end{theorem}

\begin{proof}
  \leanok
  Membership in $\cNP$ takes the set $S$ as certificate and computes the energy through the
  identity $H = -\sum J + 2\,\mathrm{cut}_J$ in polynomial time. Hardness composes the
  completeness of $\MAXCUT$ (Theorem~\ref{thm:maxcut}) with the reduction that sends a weight $w_e$
  to the coupling $J_e = -w_e$ and the cut bound $k$ to the energy bound $-\sum J - 2k$.
\end{proof}

\begin{theorem}[$\cP = \cNP$ if and only if $\ISING \in \cP$]
  \label{thm:physics-p-np-ising}
  \lean{Sagitta.Map.Informacion.Hamiltoniano.P_eq_NP_iff_ISING_mem_P}
  \leanok
  \uses{thm:physics-ising-npcomplete, def:P, def:NP}
  The following are equivalent. First, for every type $\alpha$ and every injective encoding
  $\enc : \alpha \to \{0,1\}^{*}$, the class $\cP$ for $\enc$ equals the class $\cNP$ for
  $\enc$. Second, $\ISING$ is in the class $\cP$ for the encoding of
  Definition~\ref{def:physics-ising-problem}.
  \own{}
  \leanfile{Sagitta/Map/Informacion/Hamiltoniano.lean}
\end{theorem}

\begin{proof}
  \leanok
  From the first to the second, rewrite the class $\cP$ of the encoding of $\ISING$ as $\cNP$ and
  use $\ISING \in \cNP$. From the second to the first, an $\cNP$-complete problem in $\cP$ gives
  $\cP = \cNP$ for every encoding, by Theorem~\ref{thm:physics-ising-npcomplete}.
\end{proof}

\section{Transfer along an isomorphism of structures}
\label{sec:physics-transfer}

A \emph{structure} is a group acting on a set of states together with a real energy on the states.
An isomorphism of structures carries what is defined from the group action and the energy.
The statements below are about this abstract structure and about the Ising model on $n$ sites with
a real coupling matrix; they contain no machine and no measure of cost.

\begin{definition}[Isomorphism of structures: a group acting on states with an energy]
  \label{def:physics-iso-estructura}
  \lean{Sagitta.Map.Informacion.Transferencia.IsoEstructura, Sagitta.Map.Informacion.Transferencia.EsMinimo}
  \leanok
  Let $G$ act on $X$ with energy $E : X \to \mathbb{R}$, and $G'$ act on $X'$ with energy
  $E' : X' \to \mathbb{R}$. An \emph{isomorphism of structures} consists of a bijection
  $\varphi : X \to X'$, a group isomorphism $h : G \to G'$ with $\varphi(g \cdot x) = h(g) \cdot
  \varphi(x)$ for all $g, x$, and real numbers $a > 0$ and $b$ with $E'(\varphi(x)) = a\,E(x) + b$
  for all $x$. A state $x$ is a \emph{minimum} of $E$ when $E(x) \le E(y)$ for every $y \in X$.
  \own{}
  \leanfile{Sagitta/Map/Informacion/Transferencia.lean}
\end{definition}

\begin{definition}[The Ising energy and the cut weight, on $n$ sites with a coupling matrix]
  \label{def:physics-ising-matrix}
  \lean{Sagitta.Map.Informacion.Transferencia.ising, Sagitta.Map.Informacion.Transferencia.pesoTotal, Sagitta.Map.Informacion.Transferencia.Corte, Sagitta.Map.Informacion.Transferencia.corte, Sagitta.Map.Informacion.Transferencia.ladoMas}
  \leanok
  Let $J : \{1..n\} \times \{1..n\} \to \mathbb{R}$ and let a spin configuration be
  $\sigma : \{1..n\} \to \{+1, -1\}$ (the units of $\mathbb{Z}$). The energy is
  $H_J(\sigma) = -\sum_{i<j} J_{ij}\,\sigma_i\sigma_j$ and the total weight is
  $W(J) = \sum_{i<j} J_{ij}$. A cut is a subset $S$ of the sites; its weight is
  $\mathrm{cut}_J(S) = \sum_{i<j,\ i \text{ and } j \text{ on different sides of } S} J_{ij}$. The
  side $\mathrm{ladoMas}(\sigma)$ of a configuration is the set of sites with $\sigma_i = +1$.
  \own{}
  \leanfile{Sagitta/Map/Informacion/Transferencia.lean}
\end{definition}

\begin{theorem}[Ising energy and cut weight, on $n$ sites with a real coupling matrix]
  \label{thm:physics-ising-eq-corte}
  \lean{Sagitta.Map.Informacion.Transferencia.ising_eq_corte}
  \leanok
  \uses{def:physics-ising-matrix}
  For every $n$, every $J : \{1..n\} \times \{1..n\} \to \mathbb{R}$ and every configuration
  $\sigma : \{1..n\} \to \{+1, -1\}$,
  \[
    H_J(\sigma) = -W(J) + 2\,\mathrm{cut}_J(\mathrm{ladoMas}(\sigma)).
  \]
  No symmetry of $J$ and no condition on its diagonal is assumed, because the sums run over the
  pairs $i < j$.
  \own{}
  \leanfile{Sagitta/Map/Informacion/Transferencia.lean}
\end{theorem}

\begin{proof}
  \leanok
  Each pair contributes $J_{ij}\sigma_i\sigma_j = J_{ij} - 2\,[\sigma_i \ne \sigma_j]\,J_{ij}$,
  checked on the four sign cases; sum over the pairs.
\end{proof}

\begin{theorem}[Isomorphisms of structures carry the set of minima]
  \label{thm:physics-image-minimos}
  \lean{Sagitta.Map.Informacion.Transferencia.IsoEstructura.image_minimos}
  \leanok
  \uses{def:physics-iso-estructura}
  Let $I$ be an isomorphism of structures between $(G \curvearrowright X, E)$ and
  $(G' \curvearrowright X', E')$, with bijection $\varphi$. Then the image under $\varphi$ of the
  set of minima of $E$ is the set of minima of $E'$.
  \own{}
  \leanfile{Sagitta/Map/Informacion/Transferencia.lean}
\end{theorem}

\begin{proof}
  \leanok
  A state $x$ is a minimum of $E$ if and only if $\varphi(x)$ is a minimum of $E'$: apply
  $E'(\varphi(y)) = a\,E(y) + b$ with $a > 0$ and use that $\varphi$ is onto. The two inclusions
  follow, taking $\varphi^{-1}$ for the one that needs a preimage.
\end{proof}

\begin{theorem}[Isomorphisms of structures carry stabilizers]
  \label{thm:physics-map-stabilizer}
  \lean{Sagitta.Map.Informacion.Transferencia.IsoEstructura.map_stabilizer}
  \leanok
  \uses{def:physics-iso-estructura}
  Let $I$ be an isomorphism of structures between $(G \curvearrowright X, E)$ and
  $(G' \curvearrowright X', E')$, with bijection $\varphi$ and group isomorphism $h$. For every
  $x \in X$, the image under $h$ of the stabilizer of $x$ in $G$ is the stabilizer of $\varphi(x)$
  in $G'$.
  \own{}
  \leanfile{Sagitta/Map/Informacion/Transferencia.lean}
\end{theorem}

\begin{proof}
  \leanok
  By equivariance, $g$ fixes $x$ if and only if $h(g)$ fixes $\varphi(x)$, using that $\varphi$ is
  injective. For the inclusion from right to left, a stabilizing element of $G'$ is the image of
  its preimage under $h$.
\end{proof}

\begin{theorem}[Orbit counts of $\{\pm 1\}$ on truth assignments and on spin configurations agree]
  \label{thm:physics-orbitas-sat-espin}
  \lean{Sagitta.Map.Informacion.Transferencia.orbitas_sat_espin}
  \leanok
  \uses{def:physics-iso-estructura}
  Let $n \in \mathbb{N}$ and let the group $\mathbb{Z}^{\times} = \{+1, -1\}$ act on the set of maps
  $\{1..n\} \to \{\mathrm{true}, \mathrm{false}\}$ by negating every value when the element is $-1$,
  and on the set of maps $\{1..n\} \to \{+1, -1\}$ by multiplying every value. Then the two
  actions have the same number of orbits.
  \own{}
  \leanfile{Sagitta/Map/Informacion/Transferencia.lean}
\end{theorem}

\begin{proof}
  \leanok
  The bijection from truth values to spins, coordinate by coordinate, together with the identity on
  $\{+1,-1\}$, is an isomorphism of structures for the empty formula; the orbit counts of
  an isomorphism of structures agree, by a bijection between the two quotients.
\end{proof}

\section{Random 3-SAT as a finite probability space}
\label{sec:physics-random-sat}

The statements of this section are about finite probability: a finite set $\Omega$ with weights
that sum to $1$, no measure theory. A random formula is a uniform choice of $m$ clauses, each a
triple of variables and three signs, over a finite type of triples.

\begin{definition}[A finite probability law, and the number of solutions of a clause list]
  \label{def:physics-ley}
  \lean{Sagitta.Map.FisicaEstadistica.Ley, Sagitta.Map.FisicaEstadistica.Ley.esp, Sagitta.Map.FisicaEstadistica.Ley.prob, Sagitta.Map.FisicaEstadistica.uniforme, Sagitta.Map.FisicaEstadistica.Clausula, Sagitta.Map.FisicaEstadistica.Satisface, Sagitta.Map.FisicaEstadistica.numSol, Sagitta.Map.FisicaEstadistica.probSat}
  \leanok
  A \emph{law} $P$ on a finite type $\Omega$ is a family of weights $w(\omega) \ge 0$ with
  $\sum_\omega w(\omega) = 1$; $E_P[X] = \sum_\omega w(\omega) X(\omega)$ and
  $P(A) = \sum_{\omega \in A} w(\omega)$. The uniform law gives every point the weight
  $1/|\Omega|$. Let $V$ be a finite type with a map $\mathrm{var} : V \to \{1,2,3\} \to \{1..n\}$. A
  \emph{clause} is a pair $(v, s)$ with $v \in V$ and $s : \{1,2,3\} \to \{\mathrm{true},
  \mathrm{false}\}$; an assignment $x : \{1..n\} \to \{\mathrm{true}, \mathrm{false}\}$
  \emph{satisfies} it when $x(\mathrm{var}(v)(i)) = s(i)$ for some $i$. For a formula
  $F : \{1..m\} \to \text{clauses}$, $\mathrm{numSol}(F)$ is the number of assignments that satisfy
  every clause of $F$. The probability $\mathrm{probSat}$ is the probability, under the uniform law
  on $\{1..m\} \to \text{clauses}$, that $\mathrm{numSol}(F) \ge 1$, and is $0$ when $V$ is empty.
  \own{}
  \leanfile{Sagitta/Map/Informacion/FisicaEstadistica.lean}
\end{definition}

\begin{theorem}[First moment of the number of solutions, for $m$ independent clauses over a finite type of triples]
  \label{thm:physics-esperanza-soluciones}
  \lean{Sagitta.Map.FisicaEstadistica.esperanza_soluciones}
  \leanok
  \uses{def:physics-ley}
  Let $V$ be a finite nonempty type, $\mathrm{var} : V \to \{1,2,3\} \to \{1..n\}$ and
  $m \in \mathbb{N}$. Under the uniform law on formulas of $m$ clauses, the expectation of the
  number of solutions is
  \[
    E[\mathrm{numSol}] = 2^{n}\,(7/8)^{m}.
  \]
  The triple of variables of a clause need not have distinct entries. Independence of the clauses
  is the product structure of the uniform law on $\{1..m\} \to \text{clauses}$.
  \source{\cite{FrancoPaull83}.}
  \fidelity{Not checked against the printed source.}
  \leanfile{Sagitta/Map/Informacion/FisicaEstadistica.lean}
\end{theorem}

\begin{proof}
  \leanok
  For a fixed assignment, exactly one of the $8$ sign patterns falsifies a clause, so $7/8$ of the
  clauses are satisfied. Write the number of solutions as a sum over assignments of a product over
  the $m$ clauses, exchange the sums, and factor the product.
\end{proof}

\begin{theorem}[The probability of satisfiability tends to $0$ above the first-moment threshold, for any finite types of triples]
  \label{thm:physics-prob-sat-tiende-cero}
  \lean{Sagitta.Map.FisicaEstadistica.prob_sat_tiende_cero}
  \leanok
  \uses{def:physics-ley, thm:physics-esperanza-soluciones}
  For each $n$ let $V_n$ be a finite type with $\mathrm{var}_n : V_n \to \{1,2,3\} \to \{1..n\}$,
  and assume that $V_n$ is nonempty for all large $n$. Let $c$ be a real number with
  $c > \ln 2 / \ln (8/7)$, and let $m : \mathbb{N} \to \mathbb{N}$ satisfy $c\,n \le m(n)$ for
  every $n$. Then $\mathrm{probSat}(V_n, \mathrm{var}_n, m(n))$ tends to $0$ as $n \to \infty$.
  \source{\cite{FrancoPaull83}.}
  \fidelity{Not checked against the printed source.}
  \leanfile{Sagitta/Map/Informacion/FisicaEstadistica.lean}
\end{theorem}

\begin{proof}
  \leanok
  Markov's inequality for the number of solutions and Theorem~\ref{thm:physics-esperanza-soluciones}
  bound the probability by $(2\,(7/8)^c)^n$. The ratio $2\,(7/8)^c$ is below $1$ exactly because
  $c > \ln 2 / \ln (8/7)$, so the bound tends to $0$ and the squeeze gives the limit.
\end{proof}

\begin{theorem}[Annealed entropy bounds quenched entropy, on a finite probability space]
  \label{thm:physics-templada-recocida}
  \lean{Sagitta.Map.FisicaEstadistica.templada_le_recocida}
  \leanok
  \uses{def:physics-ley}
  Let $P$ be a law on a finite type $\Omega$ and $Z : \Omega \to \mathbb{R}$ with $Z(\omega) > 0$
  for every $\omega$. Then $E_P[\log Z] \le \log E_P[Z]$.
  \source{\cite{MezardParisiZecchina02}.}
  \fidelity{Not checked against the printed source.}
  \leanfile{Sagitta/Map/Informacion/FisicaEstadistica.lean}
\end{theorem}

\begin{proof}
  \leanok
  Jensen's inequality for the concave function $\log$ on the positive reals, applied to the
  weights of the law.
\end{proof}

\begin{theorem}[Second-moment lower bound for the probability of being positive, on a finite probability space]
  \label{thm:physics-paley-zygmund}
  \lean{Sagitta.Map.FisicaEstadistica.paley_zygmund}
  \leanok
  \uses{def:physics-ley}
  Let $P$ be a law on a finite type $\Omega$ and $Z : \Omega \to \mathbb{R}$ with $Z(\omega) \ge 0$
  for every $\omega$ and $E_P[Z] > 0$. Then
  \[
    \frac{E_P[Z]^2}{E_P[Z^2]} \le P(Z > 0).
  \]
  \source{\cite{PaleyZygmund32}.}
  \fidelity{Not checked against the printed source.}
  \leanfile{Sagitta/Map/Informacion/FisicaEstadistica.lean}
\end{theorem}

\begin{proof}
  \leanok
  The Cauchy--Schwarz inequality in the form $(\sum r_i)^2 \le (\sum f_i)(\sum g_i)$ when
  $r_i^2 \le f_i g_i$, with $r_i = w_i Z_i$, $f_i = w_i Z_i^2$ and $g_i = w_i [Z_i > 0]$, gives
  $E[Z]^2 \le E[Z^2]\,P(Z>0)$. Then $E[Z^2] > 0$ and the claim is a division.
\end{proof}

\begin{theorem}[Values at the integers do not fix the derivative of the moment at $0$]
  \label{thm:physics-enteros-no-bastan}
  \lean{Sagitta.Map.FisicaEstadistica.enteros_no_bastan}
  \leanok
  \uses{def:physics-ley}
  Let $P$ be a law on a finite type $\Omega$ and $Z : \Omega \to \mathbb{R}$ with $Z(\omega) > 0$
  for every $\omega$, and let $M(r) = \sum_\omega w(\omega) Z(\omega)^r$ for real $r$. Then, for
  every natural number $k$, $M(k) + \sin(\pi k) = M(k)$; and the function
  $r \mapsto M(r) + \sin(\pi r)$ has derivative $E_P[\log Z] + \pi$ at $0$.
  \source{\cite{MezardParisiZecchina02}, for the replica method that this statement shows cannot be
  justified by integer values alone.}
  \fidelity{Not checked against the printed source.}
  \leanfile{Sagitta/Map/Informacion/FisicaEstadistica.lean}
\end{theorem}

\begin{proof}
  \leanok
  $\sin(\pi k) = 0$ for natural $k$. The derivative of $M$ at $0$ is $E_P[\log Z]$, by
  differentiating each term $w(\omega) Z(\omega)^r$, and the derivative of $\sin(\pi r)$ at $0$ is
  $\pi$; add them.
\end{proof}

\section{The distribution of a finite model of programs as a Gibbs mass}
\label{sec:physics-solomonoff}

The model is a finite set $P$ of programs with a length $\ell : P \to \mathbb{N}$ and a map
$U : P \to X$ to outputs. It is not a universal machine; the sum below is over a finite set.

\begin{definition}[Partition function, Gibbs mass, mass of an output and free energy, in a finite model of programs]
  \label{def:physics-gibbs}
  \lean{Sagitta.Map.Informacion.Repeticiones.Z, Sagitta.Map.Informacion.Repeticiones.masaGibbs, Sagitta.Map.Informacion.Repeticiones.m, Sagitta.Map.Informacion.Repeticiones.energiaLibre}
  \leanok
  Let $P$ be a finite type, $\ell : P \to \mathbb{N}$, $U : P \to X$ with $X$ having decidable
  equality. For a real $\beta$ and a finite subset $S$ of $P$, the Gibbs mass is
  $\sum_{p \in S} e^{-\beta\,\ell(p)}$, and the partition function $Z(\beta)$ is the Gibbs mass of
  all of $P$. The mass of an output $x$ is $m(x) = \sum_{U(p) = x} (1/2)^{\ell(p)}$. The free energy
  of $S$ at $\beta$ is $-\ln(\text{Gibbs mass of } S)/\beta$.
  \own{}
  \leanfile{Sagitta/Map/Informacion/Repeticiones.lean}
\end{definition}

\begin{theorem}[The mass of an output is a Gibbs mass at $\beta = \ln 2$, in a finite model of programs]
  \label{thm:physics-solomonoff-gibbs}
  \lean{Sagitta.Map.Informacion.Repeticiones.solomonoff_es_gibbs}
  \leanok
  \uses{def:physics-gibbs}
  Let $P$ be a finite type, $\ell : P \to \mathbb{N}$, $U : P \to X$ and $x \in X$. Then $m(x)$ is
  the Gibbs mass at $\beta = \ln 2$ of the set $\{p \mid U(p) = x\}$.
  \own{}
  \leanfile{Sagitta/Map/Informacion/Repeticiones.lean}
\end{theorem}

\begin{proof}
  \leanok
  Termwise, $e^{-\ln 2 \cdot n} = (1/2)^n$.
\end{proof}

\begin{theorem}[The free energy at $\beta = \ln 2$ is $-\log_2 m(x)$, in a finite model of programs]
  \label{thm:physics-energia-libre}
  \lean{Sagitta.Map.Informacion.Repeticiones.energia_libre_eq}
  \leanok
  \uses{def:physics-gibbs, thm:physics-solomonoff-gibbs}
  With the data of Definition~\ref{def:physics-gibbs} and $x \in X$, the free energy at
  $\beta = \ln 2$ of the set $\{p \mid U(p) = x\}$ equals $-\log_2 m(x)$.
  \own{}
  \leanfile{Sagitta/Map/Informacion/Repeticiones.lean}
\end{theorem}

\begin{proof}
  \leanok
  Replace the Gibbs mass by $m(x)$ with Theorem~\ref{thm:physics-solomonoff-gibbs}, and unfold
  $\log_2$ as $\ln / \ln 2$.
\end{proof}

\begin{theorem}[Kraft: the partition function at $\beta = \ln 2$ is at most $1$ for prefix-free codes, in a finite model of programs]
  \label{thm:physics-kraft-z}
  \lean{Sagitta.Map.Informacion.Repeticiones.kraft_Z}
  \leanok
  \uses{def:physics-gibbs}
  Let $P$ be a finite type, $\ell : P \to \mathbb{N}$ and $\mathrm{codigo} : P \to \{0,1\}^{*}$ an
  injective map whose range is a prefix-free set, with $\ell(p)$ the length of
  $\mathrm{codigo}(p)$ for every $p$. Then $Z(\ln 2) \le 1$.
  \own{}
  \leanfile{Sagitta/Map/Informacion/Repeticiones.lean}
\end{theorem}

\begin{proof}
  \leanok
  $Z(\ln 2) = \sum_p (1/2)^{\ell(p)}$, and by injectivity this is the Kraft sum over the finite
  prefix-free set of codewords, which is at most $1$ by the Kraft inequality in Mathlib.
\end{proof}

\section{Poincar\'e recurrence on a finite set}
\label{sec:physics-poincare}

The statements are about an injective map of a finite set into itself, iterated; there is no
measure and no phase space.

\begin{theorem}[Recurrence of an injective map on a finite set, within the cardinality]
  \label{thm:physics-poincare}
  \lean{Sagitta.Map.Informacion.Repeticiones.poincare}
  \leanok
  Let $\sigma$ be a finite type, $f : \sigma \to \sigma$ injective and $x \in \sigma$. Then there
  is $n$ with $0 < n$, $n \le |\sigma|$ and $f^{n}(x) = x$.
  \source{\cite{Poincare1890}, for the recurrence theorem; this is its discrete form.}
  \fidelity{Not checked against the printed source.}
  \leanfile{Sagitta/Map/Informacion/Repeticiones.lean}
\end{theorem}

\begin{proof}
  \leanok
  The pigeonhole principle gives $i < j \le |\sigma|$ with $f^{i}(x) = f^{j}(x)$. Since
  $f^{i}$ is injective, $f^{i}(f^{j-i}(x)) = f^{i}(x)$ gives $f^{j-i}(x) = x$.
\end{proof}

\begin{theorem}[An injective map on a finite set returns infinitely often]
  \label{thm:physics-poincare-infinitas}
  \lean{Sagitta.Map.Informacion.Repeticiones.poincare_infinitas}
  \leanok
  \uses{thm:physics-poincare}
  Let $\sigma$ be a finite type, $f : \sigma \to \sigma$ injective, $x \in \sigma$ and
  $N \in \mathbb{N}$. Then there is $n \ge N$ with $f^{n}(x) = x$.
  \source{\cite{Poincare1890}, for the recurrence theorem; this is its discrete form.}
  \fidelity{Not checked against the printed source.}
  \leanfile{Sagitta/Map/Informacion/Repeticiones.lean}
\end{theorem}

\begin{proof}
  \leanok
  Take a return time $n > 0$ from Theorem~\ref{thm:physics-poincare}. Then $x$ is periodic of period
  $n$, so it returns at $nN$, and $nN \ge N$.
\end{proof}
